\section{Del Dini derivativation}

In the main result we will need to consider the del Dini derivation of $\sup_{0 \leq t \leq T} \norm{\sol_t^N-\mfsol_t^N}_\infty$.
In this section we show how to work with such quantity.

\subsection{Del Dini derivatives and FTC}
[[define and present the del dini derivative]]

\subsection{Advance beyond the $\sup$}

In this part we want to show that

\begin{equation}
\deldiniu_T^+ \sup_{0 \leq t \leq T} \norm{\sol_t^N-\mfsol_t^N}_\infty
\leq
\sup_{0 \leq t \leq T} \deldiniu^+_t \norm{\sol_t^N-\mfsol_t^N}_\infty
\end{equation}\label{ch4:result:deldini-of-sup}

Consider the random function

\begin{equation}
\Ch(t) := | \sol_t^N - \mfsol_t^N|_\infty
\end{equation}

and look at its paths.
Such function is surely continuous between birth times $\btime_m$, but what happens at such instants?

\begin{lemma}
The paths of $\Ch(t)$ are \cadlag and have only a finite number of discontinuouties on $[0,T]$ (under $ \PrbCondTimes{}{}$).
Furthermore, considering $\Ch(t-) := \lim_{h\rightarrow 0_-} \Ch(t+h)$, we have that
\begin{equation}
\forall t_i. \quad \text{$\Ch$ is not $\Cspace^0$ in $t_i$} \Rightarrow \Ch(t_i-) \leq \Ch(t_i)
\end{equation}
\end{lemma}

\begin{remark}
Notice that this does \emph{not} necessarily mean that $\Ch$ has increasing paths, it only means that such paths, if they do have a jump, it's necessarily a jump up.
\end{remark}

\begin{figure}[h]
\begin{tikzpicture}
\end{tikzpicture}
\caption{[[[HERE DRAW TIKZ FIGURE WITH a lack of continuouty AT THE RANDOM TIMES with the jump up]]]}
\end{figure}
\label{ch4:fig:discontinuous distance}

\begin{proof}
[[write down the proof]] and quote also \ref{ch3:rmk:mf-not-c0}
\end{proof}

\begin{definition}
Let $\Skoro$ be the space of functions $g: \R_+ \longrightarrow \R_+$ such that:
\begin{itemize}
\item $g(0)=0$
\item $g$ is \cadlag
\item $g$ has finite amount of discontinuouty in $[0,T]$ for each $T$
\item If $g$ is discontinuous in $t_i$ then $g(t-) \leq g(t)$
\end{itemize}
\end{definition}

By the previous lemma the paths of $\Ch$ are in $\Skoro$.

\begin{lemma}
Let $g \in \Skoro$.
Then we have 
\begin{equation}
\sup_{0\leq s \leq t} g(s) = \max_{0\leq s \leq t} g(s)
\end{equation}
\end{lemma}

\begin{proof}
[[write down the proof]]
\end{proof}

\begin{lemma}
Let $g$ be a function in $\Skoro$.
[[require maybe finite del dini derivatives on the left side]]
Then
\begin{equation}
\deldiniu^+_t \sup_{0\leq s \leq t} g(s) 
\leq
\sup_{0\leq s \leq t} \deldiniu^+_s g(s)
\end{equation}
\end{lemma}

\begin{proof}
[[write down the proof]]
\end{proof}

Thus we proved \ref{ch4:result:deldini-of-sup}.

\subsection{Get rid of the integral sign in the derivation}

Notice that since the mean-field particles are coupled with the system \ref{ch2:model}, and since the initial conditions $\adens_0$ are the same for all the amoebas, we have that

\begin{align*}
\norm{\sol_t^N-\mfsol_t^N}_\infty 
&= \norm{
	\sol_0^N 
	+ \int_0^t \drift(s, \sol_s^N) d s
	+ 
	\sigma B_t^N
-
\left(
	\mfsol_0^N 
	+ \int_0^t \drift(s, \mfsol_s^N) d s
	+ 
	\sigma B_t^N
\right)
}
\\
&=
\norm{
	\int_0^t \drift(s, \sol_s^N) - \drift(s, \mfsol_s^N) d s
}
\end{align*}

\section{The two laws of large numbers}